% Huber & Pollmann, J. Cheminformatics 2026, 18:83, Figure 2 -- TRANSCRIBED BY HAND.
% The paper publishes these as labelled bars and gives no table, so the counts below are read
% off the printed bar labels at 400 dpi. They are their numbers, not estimates from bar lengths,
% but three were misread at screen resolution on a first pass (Morgan-3 binary 84,864 as 64,864;
% Morgan-2 binary 98,267 as 96,267; Avalon count 89,831 as 89,931) and only caught on re-render.
% Check them before publication. Percentages are computed from the counts, not transcribed.
%
% NOT the table of our own result -- that is paper/tab_degeneracy.tex, generated by
% research/fp_degeneracy_table.py, where every row is recomputed on our own compounds.
\begin{table}[t]
\centering
\begin{tabular}{lrrrr}
\toprule
& \multicolumn{2}{c}{count variant} & \multicolumn{2}{c}{binary variant} \\
\cmidrule(lr){2-3}\cmidrule(lr){4-5}
fingerprint & $n$ & \% of set & $n$ & \% of set \\
\midrule
\textbf{MAP4} & \textbf{223} & \textbf{0.03}\% & \textbf{253} & \textbf{0.04}\% \\
Atom pair & 7{,}007 & 0.98\% & 7{,}634 & 1.06\% \\
Morgan-9 & 21{,}258 & 2.96\% & 30{,}616 & 4.26\% \\
Morgan-3 & 64{,}899 & 9.04\% & 84{,}864 & 11.8\% \\
FCFP9 & 71{,}896 & 10.0\% & 86{,}882 & 12.1\% \\
RDKit & 72{,}693 & 10.1\% & 128{,}578 & 17.9\% \\
Morgan-2 & 74{,}000 & 10.3\% & 98{,}267 & 13.7\% \\
PubChem & 82{,}630 & 11.5\% & 248{,}599 & 34.6\% \\
LINGO & 82{,}796 & 11.5\% & 138{,}002 & 19.2\% \\
Topol.\ torsion & 87{,}608 & 12.2\% & 105{,}557 & 14.7\% \\
Avalon & 89{,}831 & 12.5\% & 113{,}114 & 15.8\% \\
Klekota--Roth & 93{,}817 & 13.1\% & 164{,}942 & 23.0\% \\
MACCS & 113{,}899 & 15.9\% & 309{,}279 & 43.1\% \\
FCFP3 & 116{,}997 & 16.3\% & 137{,}564 & 19.2\% \\
\addlinespace
Element counts (11-dim) & 617{,}659 & 86.0\% & \multicolumn{2}{c}{--} \\
\bottomrule
\end{tabular}
\caption{Fingerprint degeneracy reported by Huber and Pollmann~\cite{Huber2026}: the
number of compounds, out of 718{,}067 unique biomolecular structures, that share an
identical fingerprint with at least one structurally different compound. Rows are ordered
by the count variant. Counting substructure occurrences rather than mere presence improves
specificity for every fingerprint, dramatically so for the dictionary-based ones --- MACCS
collides for 43.1\% of the set as a bit vector and 15.9\% as a count vector. A plain
11-dimensional element-count vector is included as their baseline: it collides for 86\% of
the set, yet the paper reports nearly all of those collisions fall within 1\,Da, so a high
collision rate and a harmful one are not the same thing.}
\label{tab:fp-duplicates}
\end{table}
